\documentclass[11pt]{article}
\usepackage[margin=1in]{geometry}
\usepackage{amsmath,amssymb,booktabs,longtable,tabularx,array}
\usepackage{enumitem}
\usepackage[T1]{fontenc}
\usepackage[hidelinks]{hyperref}
\setlist[itemize]{leftmargin=*}
\setlist[enumerate]{leftmargin=*}
\newcommand{\code}[1]{\texttt{\detokenize{#1}}}
\newcommand{\st}{s_t}
\newcommand{\snext}{s_{t+1}}
\newcommand{\obs}{o_t}
\newcommand{\sobs}{\sigma_{\mathrm{obs}}}
\newcommand{\ssafe}{s_{\mathrm{safe}}}
\newcommand{\E}{\mathbb{E}}
\newcommand{\Var}{\mathrm{Var}}
\newcommand{\1}{\mathbf{1}}
\sloppy


\title{Problem Setting and Scientific Design}
\author{Ecology Benchmark Documentation}
\date{July 2026}

\begin{document}
\maketitle

\section{Purpose}

This document defines the scientific problem that the benchmark instantiates.
The benchmark is a partially observed ecological control problem with continuous
latent abundance, noisy surveys, finite management actions, public logged
datasets, and evaluator-only ecological truth.

\section{POMDP}

At time $t$, the latent abundance is $\st \ge 0$. Exact zero is absorbing
extinction. The agent does not observe $\st$ directly. It receives a survey
\[
  \obs = \st \epsilon_t,\qquad
  \epsilon_t \sim \mathrm{LogNormal}(0,\sobs).
\]
When $\sobs=0$, the observation equals the state up to numerical tolerance; when
$\sobs>0$, the observation is median-unbiased but has mean
$\st\exp(\sobs^2/2)$.

The public history available to policies is the observation/action sequence plus
public control fields. The true state, hidden structural parameters, regime, and
true effective growth are evaluator-only. Logged rewards are training targets in
the public dataset, but filtering and deployed policy updates do not use them as
state observations.

\section{Dynamics Families}

The environment dispatches among four continuous map families:
\begin{itemize}
  \item Ricker;
  \item Allee-Ricker;
  \item theta-logistic;
  \item regime-switching Allee-like dynamics.
\end{itemize}

For real and dummy ecology cells, management first applies direct
translocation, if any, then advances public controls. The public control
variables are
\[
  \rho_t,\quad \kappa_t,\quad
  K^{\mathrm{eff}}_t =
    \mathrm{clip}(K_{\mathrm{base}}+\kappa_t,K_{\min},K_{\max}).
\]
In set-point ecology cells, $\rho_{t+1}$ equals the chosen action's growth
set-point and $\kappa$ accumulates carrying-capacity increments. Effective growth
is
\[
  r^{\mathrm{eff}}_{t+1}
  =
  \mathrm{clip}(\rho_{t+1}, r_{\min}, r_{\max}).
\]
The code splits negative effective growth into unconditional mortality:
\[
  r_+ = \max(r^{\mathrm{eff}},0),\qquad
  r_- = \min(r^{\mathrm{eff}},0).
\]
This is a benchmark convention that avoids sign pathologies when a negative
set-point is inserted inside a density-dependent Ricker exponent.

\section{Modes}

\begin{center}
\begin{tabularx}{\linewidth}{llll}
\toprule
Mode & Source & Valid control & Purpose \\
\midrule
real & six CSVs in \code{real_ecology_data} & \code{setpoint_cumulative}
  & paper-facing ecology setting \\
dummy & \code{dummydata.py} & \code{setpoint_cumulative}
  & quick real-like experiments \\
synthetic & in-code simulators & \code{tier2_one_step}, \code{cumulative_capped}
  & stress/reference experiments \\
\bottomrule
\end{tabularx}
\end{center}

The historically confusing point is that \code{setpoint_cumulative} is shared by
real and dummy ecology. The data source is controlled by \code{data_mode}, not by
the control-mode name.

\section{Dataset Boundary}

The public dataset contains observations, actions, rewards, next observations,
done flags, episode IDs, timesteps, and public controls where enabled. The
private sidecar contains true states, next states, hidden structural parameters,
regime values, true rewards, and true effective rates. Training code accepts the
public dataset and belief caches, not the private sidecar.

This boundary is central to the benchmark. The reward is a monotone function of
the true next abundance after known costs are added back, so treating reward as a
filter observation would leak hidden state. The implementation records reward as
a supervised target and evaluation metric while keeping belief updates tied to
observations and public controls.

\section{Scientific Design Rationale}

\paragraph{Continuous state.}
The real ecology tables describe abundance and carrying capacity on physical
scales. Keeping $\st$ continuous avoids finite abundance bins and lets collapse,
MVP occupancy, and recovery be reported in the same units as the data.

\paragraph{Set-point growth.}
The real action table provides action-specific growth rates derived from
population/action $\lambda$ values. Treating the action as a set-point for $r$
keeps the intended intervention semantics: a management action selects the
current demographic regime instead of adding an arbitrary increment to a hidden
baseline.

\paragraph{Cumulative carrying capacity.}
Habitat and integrated conservation actions have lasting capacity effects. The
public accumulator $\kappa$ records those effects, and the current capacity is
clipped between the published baseline and the configured landscape ceiling.

\paragraph{Occupancy safety penalty.}
Several real populations begin depleted. A downward-crossing-only penalty would
charge a population that crosses into the unsafe region but would not charge a
population that starts and remains unsafe. The real safe setting therefore uses
occupancy by default: every step with $\snext \le \ssafe$ is penalized.

\paragraph{P-safe penalty.}
The recent real-ecology experiment selected \code{collapse_penalty = 5} for the
safe reward setting. The selection rule used the recoverable decision
populations: choose the largest tested penalty that keeps their mean safe return
non-negative while reducing collapse entries relative to $P=2$. At $P=10$, the
Puerto Rican parrot mean return is already negative; at $P=20$, both recoverable
decision populations are over-penalized. The package dataclass default is a
general fallback; experiment configs and result bundles define the penalty for a
particular run.

\section{Evaluation Setting}

Evaluation resets real and dummy ecology cells at the published or in-code
$N_0$. The collector can sample lower starts for coverage, but this does not
change gate/evaluation reset semantics. Reported metrics include operational
return, true return, unsafe occupancy, collapse entry, MVP occupancy and breach,
minimum and final abundance, persistence, economic cost, filter error, fallback
counts, and timing.

\section*{Verification}

Verified against: \path{src/real_ecology_benchmark/config.py}, \path{src/real_ecology_benchmark/envs.py}, \path{src/real_ecology_benchmark/controls.py}, \path{src/real_ecology_benchmark/dataset.py}, \path{src/real_ecology_benchmark/observation.py}, \path{src/real_ecology_benchmark/reward.py}, \path{real_ecology_runs/psafe_overnight_20260705/analysis/DECISION_psafe.md} @ be96c36

\end{document}
